%! Function Composition and Decomposition — Unit 1, Lesson 1.3 — Guided Notes
% THE in-class document, in the MAIN IDEAS / NOTES shape (LESSON_SHAPE.md §1): the vocabulary
% box, the hook, then ONE two-column guidednotes table. Each row is one idea — a label on the
% left; on the right one or two COMPLETE printed sentences, ONE large pre-drawn display, then a
% two-across grid of numbered problems with real writing room. Problems run 1..18.
% DENSITY: a \blank{} only where a number IS the answer (the handoff row of the table); no
% mid-sentence blanks; two problems across; 2–2.6cm of answer space each.
%   I DO   : the teacher reads the definition, marks up the display, works the first problem
%            of each grid (1, 3, 7, 11).
%   WE DO  : the class works the rest of each grid, students holding the pen, every one
%            cold-called.
%   YOU DO : the last row — problems 15–18, alone, for 14 minutes while the teacher circulates.
% The crux is row 4, TWO HANDOFFS: problem 12 runs ONE input through BOTH orders off the same
% pair of functions and gets two different answers. Problem 15 transfers it to a table.
% NOTE ON ORDERING: transformations are now Lesson 1.5, AFTER this lesson — nothing here may
% treat f(x+h) as prior knowledge. Inverse functions are Lesson 1.4, the next lesson.
%
% PAGE LOCKSTEP with notes_key/: the vocabulary write lines -> \ansline, \blank -> \ans, and
% the answer argument of every \pcell, \writespace and \labelbox filled. Nothing else differs.
% Inside a cell `\\` ENDS THE ROW — break lines with \par. (`\\` inside cases/tabular is safe:
% those environments scope it.) A row cannot break across pages: the figure and EACH grid row
% are separate sub-rows (`\\` then `& ...`, probgrid*).
\documentclass[10pt]{article}
\usepackage{precalculus-article}
\usepackage{precalculus-boxes}
\newcommand{\TallMath}[1]{$\displaystyle #1\rule[-1.4em]{0pt}{3.2em}$}

\begin{document}

\pageheader{Unit 1, Lesson 1.3}{Guided Notes}

\vspace{0.12in}

\boxguard
\begin{vocabbox}
Fill in each term as we define it together.
\par\vspace{2pt}
\noindent\textbf{\textcolor{plum}{Composition:}}\\[1pt]
\writeline\\[3pt]
\noindent\textbf{\textcolor{plum}{Inner and outer function:}}\\[1pt]
\writeline\\[3pt]
\noindent\textbf{\textcolor{plum}{Handoff value:}}\\[1pt]
\writeline\\[3pt]
\noindent\textbf{\textcolor{plum}{Domain of a composite:}}\\[1pt]
\writeline\\[3pt]
\noindent\textbf{\textcolor{plum}{Decomposition:}}\\[1pt]
\writeline
\end{vocabbox}

\vspace{0.1in}

\boxguard
\begin{hookbox}
A \$60 jacket, a \$10 coupon, and 6\% sales tax. At register 1 the clerk scans
the coupon first and the tax second. At register 2 the clerk scans the tax first
and the coupon second. Between the two scans, each register's little screen shows
a number --- and the two screens will not agree.

Same jacket, same coupon, same tax rate. Will the two \textit{totals} agree?
Commit to an answer before you compute anything.
\writelines{2}
\end{hookbox}

\small
\begin{guidednotes}
% ── ROW 1: composition and the handoff ───────────────────────────────────────
\mainidea[The number in between]{The Handoff} &
A \textbf{composition} runs two functions back to back: the output of the first becomes the
input of the second. We write it $(f \circ g)(x) = f\big(g(x)\big)$ and read it ``$f$ of $g$ of
$x$.'' The circle is \textbf{not} multiplication, and the function written \emph{last} is the one
that runs \emph{first}. The one number sitting between the two machines is the \textbf{handoff}:
it is the inner function's \textbf{answer} and the outer function's \textbf{question} at the same
time, so find it first and write it down.

\vspace{4pt}
A cable car climbs 200 feet every minute, so after $t$ minutes its elevation is $E(t) = 200t$
feet. The air cools as you rise: at elevation $e$ feet the temperature is
$T(e) = 68 - \dfrac{e}{100}$ degrees Fahrenheit. Nothing on the mountain measures degrees against
minutes --- but chaining the two instruments does.
\\
& \begin{center}
\begin{tikzpicture}[x=1cm, y=1cm, >=stealth]
  \node[draw=plum, very thick, fill=lilac, rounded corners=3pt, minimum width=1.3cm,
        minimum height=1.0cm, font=\large] (E) at (0,0) {$E$};
  \node[draw=plum, very thick, fill=lilac, rounded corners=3pt, minimum width=1.3cm,
        minimum height=1.0cm, font=\large] (T) at (5.4,0) {$T$};
  \node[font=\small] (in) at (-2.4,0) {$t$ minutes};
  \node[font=\small] (out) at (8.0,0) {$T(E(t))$ degrees};
  \draw[->, very thick, plum] (in) -- (E);
  \draw[->, very thick, goldacc] (E) -- (T);
  \draw[->, very thick, plum] (T) -- (out);
  \node[above, font=\small] at (2.7,0.16) {$E(t)$ feet};
  \node[above, font=\scriptsize\itshape, text=slate] at (0,0.58) {runs first};
  \node[above, font=\scriptsize\itshape, text=slate] at (5.4,0.58) {runs second};
  \node[below, font=\scriptsize\bfseries, text=goldacc] at (2.7,-0.10) {THE HANDOFF};
  \node[below, font=\scriptsize, text=slate, align=center] at (2.7,-0.44)
       {$E$'s answer \textit{is} $T$'s question};
\end{tikzpicture}
\end{center}
\vspace{2pt}
The handoff is measured in \labelbox{1.8cm}{}\qquad and $E$'s answer becomes $T$'s \labelbox{2.0cm}{}
\\
& \notesprompt{Say the handoff out loud before you compute anything.}
\begin{probgrid}{|Y|Y|}
\pcell{1}{Name the handoff in $T(E(4))$, with its units. Then find $T(E(4))$.}{2.2cm}{} &
\pcell{2}{In $(T \circ E)(t)$, which function runs first? Then name one thing $(T \circ E)(t)$ does \emph{not} mean.}{2.2cm}{} \\ \hline
\end{probgrid}
\\ \hline
% ── ROW 2: evaluating a composite from a table ───────────────────────────────
\mainidea[From a table]{Finding It} &
A composition never needs a formula --- it needs two lookups, so a table is enough. Read the
inner function's value first, write that number down, and only then carry it to the outer
function. If the outer function has no value at the handoff, the composite has no value there
either: that restriction is the \textbf{domain of a composite}.
\\
& \begin{center}
\renewcommand{\arraystretch}{1.4}
\begin{tabular}{|c|c|c|c|c|c|}
\hline
$x$ & 0 & 1 & 2 & 3 & 4 \\
\hline
$f(x)$ & 3 & 4 & 0 & 2 & 1 \\
\hline
$g(x)$ & 2 & 3 & 7 & 0 & 4 \\
\hline
\end{tabular}
\end{center}
\vspace{2pt}
\textbf{Fill the handoff row all the way across before you look up a single value of $f$.}
\begin{center}
\renewcommand{\arraystretch}{1.4}
\begin{tabular}{|c|c|c|c|}
\hline
$x$ & 0 & 3 & 4 \\
\hline
handoff: $g(x)$ & \blank{1.1cm} & \blank{1.1cm} & \blank{1.1cm} \\
\hline
$f(g(x))$ & \blank{1.1cm} & \blank{1.1cm} & \blank{1.1cm} \\
\hline
\end{tabular}
\end{center}
\\
& \notesprompt{Use the table above. Write the handoff, then the answer.}
\begin{probgrid}{|Y|Y|}
\pcell{3}{Find $f(g(1))$. Name the handoff first.}{2.0cm}{} &
\pcell{4}{Find $g(f(1))$. Name the handoff first.}{2.0cm}{} \\ \hline
\end{probgrid}
\\
& \begin{probgrid}*{|Y|Y|}
\pcell{5}{Try $f(g(2))$. What is the handoff, and why does the composite have no value here?}{2.2cm}{} &
\pcell{6}{Find $f(g(3))$ and $g(f(3))$. Write both handoffs.}{2.2cm}{} \\ \hline
\end{probgrid}
\\ \hline
% ── ROW 3: building the composite from formulas ──────────────────────────────
\mainidea[A whole expression]{Building It} &
With formulas instead of a table the handoff stops being a number and becomes an
\textbf{expression} --- and nothing else changes. To build $(f \circ g)(x)$, hand $f$ the whole
rule of $g$:

\vspace{3pt}
\stepnum{1}\ Write $f$ with an empty slot: $f(\square) = 3(\square) - 2$. \quad
\stepnum{2}\ Put $g(x)$ in \emph{every} slot --- circle each $x$ first so you miss none. \quad
\stepnum{3}\ Simplify.
\\
& \begin{center}
\begin{tikzpicture}[scale=0.95, >=stealth]
  \node[font=\large] (rule) at (0,0) {$f(\square) = 3(\square) - 2$};
  \node[font=\large] (fill) at (6.3,0) {$f(x^2 + 1) = 3(x^2 + 1) - 2$};
  \draw[->, very thick, plum] (2.1,0) -- node[above, font=\scriptsize\itshape] {the handoff goes in} (3.4,0);
  \node[font=\large] at (6.3,-1.5) {$= 3x^2 + 1$};
  \node[draw=goldacc, thick, fill=hookbg, rounded corners=2pt, font=\small, align=left]
        at (0,-1.5) {\textbf{The trap:} $f(g(x))$ is \emph{not} $f(x)\cdot g(x)$.\\ A product has no handoff.};
\end{tikzpicture}
\end{center}
\vspace{2pt}
Throughout this row, $f(x) = 3x - 2$ and $g(x) = x^2 + 1$.
\\
& \notesprompt{Build each composite and simplify. Show the substitution before you expand.}
\begin{probgrid}{|Y|Y|}
\pcell{7}{Build $(f \circ g)(x)$ for $f(x) = 3x-2$ and $g(x) = x^2+1$.}{2.2cm}{} &
\pcell{8}{Now build $(g \circ f)(x)$ for the same two functions. Is it the same expression?}{2.2cm}{} \\ \hline
\end{probgrid}
\\
& \begin{probgrid}*{|Y|Y|}
\pcell{9}{Check $f(g(2))$ two ways: by naming the handoff, and from your answer to problem 7.}{2.2cm}{} &
\pcell{10}{Let $p(x) = x^2$ and $q(x) = x - 4$. Build $(p \circ q)(x)$ and $(q \circ p)(x)$.}{2.4cm}{} \\ \hline
\end{probgrid}
\\ \hline
% ── ROW 4: THE CRUX — two orders, two handoffs — plus decomposition ──────────
\mainidea[Order matters]{Two Handoffs} &
Swap the two machines and you swap the handoff --- and a different number handed over produces a
different answer. That is why $f(g(x))$ and $g(f(x))$ are usually \emph{different functions}, not
two ways of writing the same one. Back at the register, the coupon is $c(p) = p - 10$ and the tax
is $t(x) = 1.06x$.
\\
& \begin{center}
\begin{tikzpicture}[x=1cm, y=1cm, >=stealth]
  % register 1 — coupon, then tax
  \node[font=\small\bfseries, text=plum, anchor=west] at (-1.5,1.85) {Register 1 --- coupon first, then tax};
  \node[font=\small] (p1) at (-0.7,1.15) {\$60};
  \node[draw=plum, very thick, fill=lilac, rounded corners=2pt, minimum width=0.9cm,
        minimum height=0.62cm, font=\small] (c1) at (1.0,1.15) {$c$};
  \node[draw=goldacc, very thick, fill=hookbg, rounded corners=2pt, minimum width=1.3cm,
        minimum height=0.62cm, font=\small\bfseries] (h1) at (2.9,1.15) {\$50};
  \node[draw=plum, very thick, fill=lilac, rounded corners=2pt, minimum width=0.9cm,
        minimum height=0.62cm, font=\small] (t1) at (4.8,1.15) {$t$};
  \node[font=\small\bfseries] (o1) at (6.7,1.15) {\$53.00};
  \draw[->, very thick, plum] (p1) -- (c1);
  \draw[->, very thick, goldacc] (c1) -- (h1);
  \draw[->, very thick, goldacc] (h1) -- (t1);
  \draw[->, very thick, plum] (t1) -- (o1);
  % register 2 — tax, then coupon
  \node[font=\small\bfseries, text=plum, anchor=west] at (-1.5,0.15) {Register 2 --- tax first, then coupon};
  \node[font=\small] (p2) at (-0.7,-0.55) {\$60};
  \node[draw=plum, very thick, fill=lilac, rounded corners=2pt, minimum width=0.9cm,
        minimum height=0.62cm, font=\small] (t2) at (1.0,-0.55) {$t$};
  \node[draw=goldacc, very thick, fill=hookbg, rounded corners=2pt, minimum width=1.3cm,
        minimum height=0.62cm, font=\small\bfseries] (h2) at (2.9,-0.55) {\$63.60};
  \node[draw=plum, very thick, fill=lilac, rounded corners=2pt, minimum width=0.9cm,
        minimum height=0.62cm, font=\small] (c2) at (4.8,-0.55) {$c$};
  \node[font=\small\bfseries] (o2) at (6.7,-0.55) {\$53.60};
  \draw[->, very thick, plum] (p2) -- (t2);
  \draw[->, very thick, goldacc] (t2) -- (h2);
  \draw[->, very thick, goldacc] (h2) -- (c2);
  \draw[->, very thick, plum] (c2) -- (o2);
  % the handoff column
  \draw[goldacc, very thick, dashed] (2.9,1.62) -- (2.9,-1.05);
  \node[below, font=\scriptsize\bfseries, text=goldacc] at (2.9,-1.05) {the two handoffs};
\end{tikzpicture}
\end{center}
\vspace{2pt}
The two totals differ by \labelbox{1.6cm}{} because the two \labelbox{2.2cm}{} differ
\\
& \fcolorbox{plum}{lilac}{\parbox{\dimexpr\linewidth-2\fboxsep-2\fboxrule\relax}{\centering
\textbf{Swap the machines and you swap the handoff. Taxing first charges 6\% of the ten dollars
you never paid --- at every price. Same two rules, different middle number, different answer:
that is what ``order matters'' actually means.}}}

\vspace{3pt}
\textbf{In your own words:} why is it not enough to say ``$f(g(x))$ and $g(f(x))$ use the same two rules''?
\writespace{1.4cm}{}
\\
& \textbf{Backwards: decomposition.} Running the machines backwards means taking one function
apart. Ask what a calculator would compute \emph{first}: that value is the handoff, and its rule
is the \textbf{inner} function. Whatever is done to it is the \textbf{outer} function. Reading
$h(x) = \sqrt{x + 5}$ this way: you would add 5 first, so the inner rule is $g(x) = x + 5$ and
the outer rule is $f(u) = \sqrt{u}$. Check a split by substituting back --- $f(g(x))$ must
rebuild $h$ exactly.
\\
& \notesprompt{Name the handoff every time. On a decomposition, check by substituting back.}
\begin{probgrid}{|Y|Y|}
\pcell{11}{Build $t(c(p))$ and $c(t(p))$ symbolically. By how much do the two rules differ at any price?}{2.4cm}{} &
\pcell{12}{Let $f(x) = 2x$ and $g(x) = x - 6$. Find $f(g(10))$ and $g(f(10))$. Write both handoffs and both answers.}{2.4cm}{} \\ \hline
\end{probgrid}
\\
& \begin{probgrid}*{|Y|Y|}
\pcell{13}{Decompose $h(x) = \sqrt{x+5}$: give the inner rule and the outer rule, then check by substituting back.}{2.2cm}{} &
\pcell{14}{A classmate says $f(g(x))$ and $g(f(x))$ must match because they use the same two rules. Use the word \emph{handoff} to say what is wrong.}{2.2cm}{} \\ \hline
\end{probgrid}
\\ \hline
% ── LAST ROW: YOU DO — alone, while the teacher circulates (14 min) ──────────
\mainidea[You do, alone]{Table \& Formulas} &
A pair of functions $p$ and $q$ given only by a table, and a pair $f$ and $g$ given only by
formulas. Name the handoff before every answer.
\\
& \begin{center}
\begin{tikzpicture}[scale=1.0]
  \node[font=\normalsize, anchor=north] at (0,1.9) {%
    \renewcommand{\arraystretch}{1.3}%
    \begin{tabular}{|c|c|c|c|c|c|}
    \hline $x$ & 1 & 2 & 3 & 4 & 5 \\ \hline $p(x)$ & 5 & 3 & 1 & 4 & 2 \\
    \hline $q(x)$ & 4 & 1 & 2 & 5 & 3 \\ \hline
    \end{tabular}};
  \node[draw=plum, very thick, fill=lilac, rounded corners=3pt, inner sep=6pt,
        font=\normalsize, anchor=north] at (0,-0.55) {$f(x) = x + 3$ \qquad $g(x) = x^2$};
\end{tikzpicture}
\end{center}
\\
& \notesprompt{On your own. Show your work.}
\begin{probgrid}{|Y|Y|}
\pcell{15}{From the table: find $p(q(2))$ and $q(p(2))$. Write each handoff. Are the two answers equal?}{2.2cm}{} &
\pcell{16}{Using $f$ and $g$ above, build $(f \circ g)(x)$ and $(g \circ f)(x)$, simplified. Same function or not?}{2.4cm}{} \\ \hline
\end{probgrid}
\\
& \begin{probgrid}*{|Y|Y|}
\pcell{17}{Find $f(g(-2))$ by naming the handoff first, then check it with your formula from problem 16.}{2.2cm}{} &
\pcell{18}{Decompose $m(x) = (2x+1)^3$ into an inner rule and an outer rule. Which rule runs first when you compute $m(1)$? Find $m(1)$.}{2.4cm}{} \\ \hline
\end{probgrid}
\\ \hline
\end{guidednotes}

% ── THE NOTES END AT THE YOU DO ROW; AP Practice and the Homework follow ──────

\end{document}
